\chapter{模为任意值的二进制计数器}

\section{通用设计公式}

前面我们分别讨论了模12、模24、模60。现在给出\textbf{适用于任意模值N的通解}。

\begin{designflow}[模N计数器通用设计]
\textbf{步骤1：确定位宽}
$$k = \lceil \log_2 N \rceil$$

\textbf{步骤2：确定清零条件}
$$\text{清零} \iff cnt = N-1$$

\textbf{步骤3：编写VHDL}
\begin{lstlisting}
 if rising_edge(clk) then
  if cnt = (N-1) then
    cnt <= (others => '0');
  else
    cnt <= cnt + 1;
  end if;
end if;
\end{lstlisting}
\end{designflow}

\section{为什么这个模板适用于任意N}

计数器的行为是确定的：
\begin{enumerate}
  \item 状态序列：$0 \to 1 \to 2 \to \cdots \to N-1 \to 0 \to 1 \to \cdots$
  \item 转移规则：$S_{next} = S_{current} + 1$（除了在$N-1$处）
  \item 在$S = N-1$处：$S_{next} = 0$
\end{enumerate}

这个行为的硬件实现正好对应上述VHDL模板。

\section{实例：模7计数器}

\begin{enumerate}
  \item $N = 7$
  \item $k = \lceil \log_2 7 \rceil = 3$（3位，最多表示7）
  \item 清零条件：cnt = 6(110)
\end{enumerate}

状态序列：000 $\to$ 001 $\to$ 010 $\to$ 011 $\to$ 100 $\to$ 101 $\to$ 110 $\to$ 000 $\to$ ...

\begin{lstlisting}[caption={模7计数器}]
signal cnt : std_logic_vector(2 downto 0);
...
 if rising_edge(clk) then
  if cnt = "110" then    -- = 6
    cnt <= "000";
  else
    cnt <= cnt + 1;
  end if;
end if;
\end{lstlisting}

\section{实例：模5计数器}

\begin{enumerate}
  \item $N = 5$
  \item $k = \lceil \log_2 5 \rceil = 3$
  \item 清零条件：cnt = 4(100)
\end{enumerate}

状态序列：000 $\to$ 001 $\to$ 010 $\to$ 011 $\to$ 100 $\to$ 000 $\to$ ...

\section{实例：模100计数器}

\begin{enumerate}
  \item $N = 100$
  \item $k = \lceil \log_2 100 \rceil = 7$（$2^6=64<100<2^7=128$）
  \item 清零条件：cnt = 99(1100011)
\end{enumerate}

\section{边界情况分析}

\subsection{$N = 2^k$ 的情况}

当N恰好是2的幂（如模16：$N = 2^4$）：
\begin{itemize}
  \item 不需要显式清零逻辑——自然溢出即可
  \item 但加上清零逻辑也能工作（只是多余）
  \item \textbf{考试注意}：模$2^N$计数器不需要if cnt=N-1，自然溢出即可
\end{itemize}

\subsection{$N = 1$ 的情况}

模1计数器：只有状态0。每个时钟周期都是0→0→0→...
但实际中没人设计模1计数器——无意义。

\section{考试常见变式}

\begin{enumerate}
  \item \textbf{变式1：给定模值N，写完整VHDL}

  直接套用模板。唯一需要注意：位宽要正确（考$\lceil \log_2 N \rceil$）

  \item \textbf{变式2：给定计数器电路图，判断模值}

  看图识别清零逻辑：清零条件是cnt = X，则模值 = X + 1。

  \item \textbf{变式3：带使能的模N计数器}

  增加enable输入：enable=1时计数，enable=0时保持。

  \item \textbf{变式4：带加载初值的模N计数器}

  增加load输入和初值输入：load=1时cnt = 初值，load=0时正常计数。

  \item \textbf{变式5：递减计数器（倒计时）}

  Count $\leftarrow$ Count - 1，到0时加载N-1。

  \item \textbf{变式6：模N计数器级联（如模60 = 模10×模6）}

  详见后续BCD计数器章节。
\end{enumerate}

\section{统一思考题}

\begin{enumerate}
  \item 设计模13计数器（$N=13$），写出VHDL。
  \item 设计模30计数器（$N=30$），写出VHDL。
  \item 设计模9计数器（$N=9$），分析需要多少位寄存器，写出VHDL。
\end{enumerate}

\textbf{答案提示}：全部套用同一模板，只改N-1的值和位宽。

\begin{examtip}[模N计数器秒杀法]
1. $k = \lceil \log_2 N \rceil$ 确定位宽
2. 清零条件 = N-1 的二进制
3. 套用模板：\texttt{if cnt = N-1 then cnt <= 0; else cnt <= cnt+1;}

这就是所有模N计数器的答案。
\end{examtip}

\section{变式详解}
\chapter{计数器——统一框架变式详解}

\section{变式1：带使能的模N计数器}

\begin{detailbox}[完整教学]
\textbf{【电路】}在标准计数器上增加en控制。
\textbf{【VHDL】}
\begin{lstlisting}
if rising_edge(clk) then
  if en = '1' then
    if cnt = N-1 then cnt <= 0; else cnt <= cnt + 1; end if;
  end if;
end if;
\end{lstlisting}
\textbf{【考点】}en=0时不赋值→保持。这是所有可控计数器的基本范式。
\end{detailbox}


\section{变式2：带同步加载的模N计数器}

\begin{detailbox}[完整教学]
\textbf{【功能】}load=1时cnt=data\_in。load=0时正常计数。
\textbf{【优先级】}load > 计数（if-else中load在前=优先级高）。
\textbf{【VHDL】}
\begin{lstlisting}
if load = '1' then cnt <= data_in;
elsif cnt = N-1 then cnt <= 0;
else cnt <= cnt + 1; end if;
\end{lstlisting}
\end{detailbox}


\section{变式3：递增/递减双向计数器}

\begin{detailbox}[完整教学]
\textbf{【VHDL】}
\begin{lstlisting}
if up = '1' then cnt <= cnt + 1; else cnt <= cnt - 1; end if;
\end{lstlisting}
\textbf{【递减自然溢出】}$0000-1=1111$（4位无符号），自动模16循环。
\end{detailbox}


\section{变式4：倒计时计数器（预置→0→预置）}

\begin{detailbox}[完整教学]
\textbf{【VHDL】}
\begin{lstlisting}
if cnt = 0 then cnt <= PRE;  -- 到0后重新加载
else cnt <= cnt - 1; end if;
\end{lstlisting}
\end{detailbox}


\section{变式5：级联计数器——总模值=乘积}

\begin{detailbox}[完整教学]
\textbf{【级联方式】}低位carry→高位en。总M=$M_1 \times M_2$。

\textbf{【实例】}模8+模5级联→总模40。模60(6×10)BCD级联。

\textbf{【VHDL关键】}高位以低位carry为时钟使能——确保高位只在低位满一轮时才+1。
\end{detailbox}


\section{变式6：格雷码计数器}

\begin{detailbox}[完整教学]
\textbf{【特点】}相邻状态仅1位不同。适合FIFO指针、跨时钟域。
\textbf{【VHDL】}case枚举所有格雷码状态。
\textbf{【3位格雷码序列】}000-001-011-010-110-111-101-100-000。
\end{detailbox}
